\documentclass[12pt]{article} % 使用12pt字号的文章类文档
\usepackage[utf8]{inputenc} % 设置输入编码为UTF-8
\usepackage[T1]{fontenc} % 使用T1字体编码
\usepackage{graphicx} % 用于插入图形
\usepackage{xcolor} % 用于颜色支持
\usepackage[UTF8]{ctex} % 支持中文

% 导入必要的包
\usepackage{tikz} % 用于程序化创建图形
\usepackage{calc} % 在 LaTeX 中进行计算
\usepackage{booktabs} % 创建专业的表格

% 定义自定义颜色
\definecolor{color1}{HTML}{000060} % 深蓝色
%\definecolor{color1}{HTML}{8C260F} % 注释掉的备用深红色
\definecolor{color2}{HTML}{333333} % 深灰色

% 使用 fontspec 包设置字体
\usepackage{fontspec}
\defaultfontfeatures{Mapping=tex-text} % 特殊字符映射
\setmainfont % 主字体设置
[BoldFont=Lato-Bold.ttf, % 粗体字体
ItalicFont=Lato-Italic.ttf, % 斜体字体
BoldItalicFont=Lato-BoldItalic.ttf] % 粗斜体字体
{Lato-Regular.ttf} % 正文字体
\newfontfamily\headingfont[ItalicFont=Lato-BlackItalic.ttf]{Lato-Black.ttf} % 标题字体

% 页面布局设置
\usepackage{geometry}
\geometry{a4paper, % A4 纸张
	hmargin=20mm,vmargin=20mm, % 页边距设置
	head=0ex,foot=0ex} % 页眉和页脚

\linespread{1.3} % 行距设置

% 设置标题和标题格式
\usepackage[hang]{caption}
\DeclareCaptionFormat{upper}{#1#2\uppercase{#3}\par} % 设置标题格式为大写
\captionsetup{labelfont={bf,color=color2},textfont={normalsize,color=color2},format = upper,figurename=图,tablename=表}

% 自定义章节格式
\usepackage{titlesec}
%\titleformat{\chapter}{\headingfont\LARGE\bfseries\scshape\color{color1}}{\thechapter}{1em}{}[\titlerule]
\titleformat{\section}{\color{color1}\headingfont\Large\bfseries\uppercase}{\thesection}{1em}{}[\titlerule]
\titleformat{\subsection}{\color{color1}\headingfont\large\bfseries\uppercase}{\thesubsection}{1em}{}
\titleformat{\subsubsection}{\color{color1}\headingfont\bfseries}{\thesubsubsection}{1em}{}

% 设置页眉和页脚
\usepackage{fancyhdr}
\pagestyle{fancy}
\lhead{} % 左页眉
\chead{} % 中间页眉
\makeatletter
\rhead{\color{color2}\@date} % 右页眉显示日期
\makeatother
\newlength{\myheight}
\lfoot{
	\settoheight{\myheight}{\thepage}
	\raisebox{-3ex-0.5\myheight}{\includegraphics[height=4ex]{logo.png}} % 左页脚显示 logo 图片
}
\cfoot{\color{color2}} % 中间页脚
\rfoot{\color{color2}\thepage} % 右页脚显示页码
\renewcommand\headrulewidth{0pt} % 去掉页眉下的横线
\renewcommand\footrulewidth{0pt} % 去掉页脚上的横线

% 自定义封面
\makeatletter
\newcommand*\DefVar[1]{\@namedef{#1}##1{\global\@namedef{get#1}{##1}}}
\DefVar{summary}
\renewcommand{\maketitle}{
	\vfill 
	\begin{center}
		% 插入标题
		\begin{tikzpicture}
			\node[draw=none, fill=color1, inner sep=10pt, text width=\textwidth-20pt, text centered] {\color{white}\headingfont\bfseries\huge\@title};
		\end{tikzpicture}
		
		% 去除标题和图片之间的缝隙
		\vspace{-2mm}
		
		% 插入封面图片，并调整其位置和大小
		\includegraphics[width=1\textwidth]{opening.jpg}
		
		% 插入作者和日期
		\vspace{50mm}
		\headingfont\bfseries\Large\@author
	\end{center}
	\vfill
	\clearpage
}
\makeatother

% 自定义盒子
\usepackage{tcolorbox}
\usepackage{wrapfig}
\def\fullboxbegin{
	\bigskip
	\begin{tcolorbox}[colback=color1,colframe=color1,coltext=white,arc=0mm,boxrule=0pt]
	}
	\def\fullboxend{\end{tcolorbox}\medskip}
%
\def\leftboxbegin{
	\begin{wrapfigure}{l}{0.5\textwidth}
		\begin{tcolorbox}[colback=color1,colframe=color1,coltext=white,arc=0mm,boxrule=0pt]
		}
		\def\leftboxend{
		\end{tcolorbox}
	\end{wrapfigure}
}
%
\def\rightboxbegin{
	\begin{wrapfigure}{r}{0.5\textwidth}
		\begin{tcolorbox}[colback=color1,colframe=color1,coltext=white,arc=0mm,boxrule=0pt]
		}
		\def\rightboxend{
		\end{tcolorbox}
	\end{wrapfigure}
}
%
\newcounter{frames}
\def\frameboxbegin#1{
	\bigskip
	\refstepcounter{frames}
	\begin{tcolorbox}[colback=white,colframe=color1,arc=0mm,title={\MakeUppercase{\textbf{ \arabic{frames}}: #1}}]
	}
	\def\frameboxend{
	\end{tcolorbox}
}
 % 输入自定义的defs文件
\usepackage{hyperref}
\usepackage{lipsum} % 用于生成示例文本

%%%%%%%%%%%%%%%
% 标题页
\title{河北建筑工程学院新生指南（2025）}

\author{本文档来源\href{https://tieba.baidu.com/p/9062077312?pn=1}{河北建筑工程学院贴吧}，共创地址在文末，\href{https://github.com/ditian-tai/HBACE-Guide-for-Freshmen}{欢迎加入一同建设}，有问题可以联系文末的学长学姐，你的提问收录会帮助更多的同学}

\date{\today} % 日期

%%%%%%%%%%%%%%%

\usepackage{fancyhdr} % 用于设置页眉和页脚
\pagestyle{fancy} % 设置页面风格为fancy
\fancyhf{} % 清空默认的页眉和页脚

% 设置页眉和页脚
\lhead{} % 左页眉
\chead{} % 中间页眉
\lfoot{
	\settoheight{\myheight}{\thepage}

}
\cfoot{} % 中间页脚为空

\renewcommand\headrulewidth{0pt} % 去掉页眉下的横线
\renewcommand\footrulewidth{0pt} % 去掉页脚上的横线

% 去掉第一页页码
\fancypagestyle{plain}{%
	\fancyhf{} % 清空页眉和页脚
	\renewcommand{\headrulewidth}{0pt} % 去掉页眉下的横线
	\renewcommand{\footrulewidth}{0pt} % 去掉页脚上的横线
}

\begin{document}
	
	\maketitle % 生成标题页
	\thispagestyle{empty} % 设置标题页无页眉和页脚
	
	\tableofcontents % 生成目录
	\clearpage
	
	% 恢复普通页面样式（无日期）
	\pagestyle{fancy}
	\fancyhf{} % 清空默认的页眉和页脚
	\lfoot{
		\settoheight{\myheight}{\thepage}
		\raisebox{-3ex-0.5\myheight}{\includegraphics[height=5ex]{logo.png}} % 左页脚显示 logo 图片
	}
	\cfoot{} % 中间页脚为空
\rfoot{\color{color2}\thepage} 
	
\section{“写在河北建筑大学前”}
	\subsection*{本文档更新于2025/6/25}

目前已知学校更名大学列入河北省“十四五”高等学校设置规划，但并不意味“十四五”一定更名成功。\\

在我看来，早晚一定能更名成功（不过可能有点晚）。\\

	\begin{figure}[!h]
	\centering
	\includegraphics[width=0.5\textwidth]{F:/latex/zhinan/fig1.jpg}
	\caption{改名体验卡} % 图注
	\end{figure}
	
	\newpage
	
%	\subsection{Boxes}
	
%	\begin{verbatim}
%		\fullboxbegin
%		Content
%		\fullboxend
%	\end{verbatim}
	
%	\begin{verbatim}
%		\leftboxbegin
%		Content
%		\leftboxend
%	\end{verbatim}
	
%	\begin{verbatim}
%		\rightboxbegin
	%	Content
%		\rightboxend
%	\end{verbatim}
	
	%\begin{verbatim}
%		\frameboxbegin{Frame Title}
	%	Content
%		\frameboxend
%	\end{verbatim}
	
%	\newpage
	
\section{基本环境介绍} 
	\subsection{气候}
	
	张家口很干，要多喝水。可以适当涂点保湿乳。张家口的风很大，天气多变。
	
		两张比较有代表性的天气图：
	\begin{figure}[!h]
		\centering
		\includegraphics[width=0.8\textwidth]{20256252.jpg}
		\caption{沙尘暴} % 图注
	\end{figure}
	\begin{figure}[!h]
		\centering
		\includegraphics[width=0.8\textwidth]{2025625.jpg}
		\caption{雷阵雨} % 图注
	\end{figure}
	
	\begin{itemize}
	\item 夏天：天气比较热，做好防晒，女孩子可以带把伞，出门用来遮阳。军训要做好防晒，以防晒伤。建议准备风扇，鉴于宿舍有时断电，最好带蓄电多的，并且充电宝能用的那种，建议带两个，一中一小，比较合适，太大的怕宿舍功率带不起来。\\
	
	\item 冬天：下雪之后特别冷，气温低，雪也不好化。冬天下雪后出门一定要慢慢走，不然容易摔跤。体质不太好的要多穿衣服，最冷的那段时间尽量少出门。不过厚被子军训的时候不建议带，除了10号楼，其他的楼每间宿舍并不大，尤其是一楼没有阳台，厚被子没地方放。\\
	


	\end{itemize}
	\subsection{宿舍} % 第一小节	
	
本科生宿舍八人间，没空调，有风扇，冬暖，夏更暖。宿舍床的面积是90cm*180cm或90cm*200cm。可以挂窗帘，下铺买床帘的话建议1.5m，上铺1.3m。除一楼外每个宿舍都有一个阳台，有纱窗。每天都会有人清洁走廊厕所。\\

大部分都是每一层有一个盥洗室及厕所，一个饮水机（刷脸付款），两个吹风机（扫码付款）。每栋宿舍一楼有一个洗衣房，洗衣机洗衣服扫码付款。宿舍饮水机用“胖乖生活”APP，支付宝扫就行，之后就可以刷脸。\\

早上7点来电，9点断电，中午11点来电，下午3点断电，晚上11点后会断电，夏天会有调整。宿舍限制功率通知是800，需要看具体宿舍，但不要在宿舍用大功率的电器，容易跳闸断电，且很危险。\\

晚上11点后宿舍大门会锁，晚回来要找舍管阿姨开，是否查寝看导员。\\

在宿舍连接wifi需要冲网条，必须要用校园卡，套餐流量够用不办校园卡也行。\\

被子建议四斤的，天冷就开暖气了，暖气挺足的。\\

其中十号楼的宿舍面积最大。所有宿舍楼一楼没阳台。\\

极少部分宿舍为六人间，上床下桌(宿舍号带H)。\\

最近听说要扩建宿舍。

\begin{figure}[!h]
	\centering
	\includegraphics[width=0.8\textwidth]{20240717155725.jpg}
	\caption{大部分宿舍图} % 图注
\end{figure}
\begin{figure}[!h]
	\centering
	\includegraphics[width=0.8\textwidth]{202407251.jpg}
	\caption{十号楼一楼宿舍图} % 图注
\end{figure}

%	\fullboxbegin
%	\lipsum[1] % 生成示例文本
%	\fullboxend

	\subsection{快递} 
截止到目前，校园内快递为京东，顺丰，极兔，韵达，其余快递均在西门校园外，床垫一类的大件建议别发其他快递的，不然取的话会很麻烦。
		\begin{figure}[h]
		\centering
		\includegraphics[width=0.5\textwidth]{20240717154657.png}
		\caption{校园外快递位置示意图} % 图注
	\end{figure}
	
		%\newpage
	
	
	\subsection{澡堂} 
	澡堂进行了全面的升级（应该是全校最豪华的地方）
	
	
	
	\subsection{食堂}
	学校共三个食堂，每个三层，可以刷微信、支付宝、校园卡，价格很平价。
	
	更新：现在食堂名改为学一、学二、学三食堂，分别对应以前的荟香阁、金源、瑞贤。
	
	荟香阁三楼会开放到晚上十点多。不出售酒。\\
	
	目前饮品有蜜雪冰城，幸运咖，甜啦啦。
	
	\subsection{教学楼}
	目前学校主要教学楼两栋，明德楼和立德楼，简称主教和副教。除了外语及建艺，其他院系都处于“漂泊”的状态，没有固定的教室学习（做毕设会分配固定的教室）。\\
	
	教学楼有免费的wifi覆盖。\\
	
	主教八楼有咖啡厅。一个教室没课就可以在那学习，每年都会有一个教室的无课表。\\
	
		\begin{figure}[!h]
		\centering
		\includegraphics[width=0.8\textwidth]{20240717163012.jpg}
		\caption{明德楼图} % 图注
	\end{figure}

		
		
	\subsection{图书馆}
	图书馆开馆时间是早上7点，闭馆时间是晚上10点。\\
	
	图书馆需要刷脸进入（新生第一次进入需要进行答题，学校会有通知），不能穿拖鞋及背心进入。\\
	
	图书馆一共6层，从东门进一楼是机房以及新能源报告厅，从北门进去上台阶是二楼，一楼和二楼不连通。每年都会有很多同学第一次去图书馆，坐电梯按二楼发现没法上去。\\
	
	图书馆四楼有研修室，扫码预约可以和同学一起自习。图书馆六楼有省级众创空间，除非是项目成员否则不能随便进入。\\
	
	每个座位上的桌子上都有二维码，记得扫码预约座位，离开点击暂离，要不东西会按占座处理收走，取回来就会有几天不能入馆的惩罚。\\
	
	图书馆大部分同学都是自习的，想要阅读推荐二楼的期刊阅览室，书会定期采购更换。\\
	
	图书馆每年举办的活动很多，奖品也很多，大家可以关注。\\
	
	\begin{figure}[!h]
		\centering
		\includegraphics[width=0.8\textwidth]{202444.jpg}
		\caption{图书馆内部图} % 图注
	\end{figure}
		\begin{figure}[!h]
		\centering
		\includegraphics[width=0.5\textwidth]{20247252.jpg}
		\caption{图书馆开放空间内部图} % 图注
	\end{figure}
			\begin{figure}[!h]
		\centering
		\includegraphics[width=0.5\textwidth]{20247253.jpg}
		\caption{研修室内部图} % 图注
	\end{figure}
		\newpage
	\subsection{各院系楼}
	除了外语，经管，数理（都在行政楼），建艺(在主教)，其他院系都有自己的楼，老师的办公和学院的实验室都在自己的楼。
	
		\subsection{行政楼}
	崇德楼，又名行政楼，主要是校领导办公的地方。也有大学物理实验室及以上几个院系办公。
	\newpage
		
%	\leftboxbegin
%	Lorem ipsum dolor sit amet, consectetuer adipiscing elit. Ut purus elit, vestibulum ut, placerat ac, adipiscing vitae, felis. Curabitur dictum gravida mauris. Nam arcu libero, nonummy eget, consectetuer id, vulputate a, magna. Donec vehicula augue eu neque. 
%	\leftboxend
	
%	\lipsum[1-2] % 生成示例文本
	
%	\rightboxbegin
%	\begin{itemize}
%		\item Lorem ipsum
%		\item Lorem ipsum
%	\end{itemize}
%	\rightboxend
	
%	\lipsum[1] % 生成示例文本
	
%	\subsubsection{First subsubsection} % 第一小节的子节
	
%	\lipsum[1] % 生成示例文本
	
%	\begin{figure}[!h]
%		\centering
%		\includegraphics[width=0.5\textwidth]{20240717154657.png}
		%\caption{The sky is the limit.} % 图注
%	\end{figure}
	
	\section{入学规划} 
	\subsection{概况}
	未来大部分人都是几种未来，出国、创业、考公、考研、就业、继承家业，这几种类型我身边都有接触。\\
	
	在学校我接触过的老师都是负责任的，讲课也很认真，期末大部分都会捞（平时分给的很高），不会有什么恶意针对的情况。\\
	
	现在监考变严了，每年到了考试周会有很多作弊被抓的同学公示，不过一般期末成绩卷面分考个四五十分左右就能通过，重修的是那种卷面成绩太低，老师实在没办法给分的。
	
	\subsection{学校各组织的基本情况}
	刚入学会有军训，可以熟悉同学及学校的基本情况。(军训休息需要坐马扎，就是木头可折叠的那种，可自己带或者找学长学姐买)\\
	
	军训期间会有各类社团及组织招新，这个时候就需要提前开始考虑大学期间的打算。\\
	
	校内的组织分为三种，校级组织，院级组织，学生社团及工作室项目组。\\
	
	
	校级组织有团委，自律会，学生会，社联，教服，图志协，校青协，大记者团，大艺团。\\
	
	所有的校级组织都有自己的公众号，直接搜索河北建筑工程学院+相应的组织就可以了解其活动。有些不太了解，我就不具体介绍了。\\
	
	院系的组织主要是各个院系的学生会。（不公开评价，想加入的话多咨询加入及未加入的学长学姐再做决定。）\\
	
	学生社团及工作室项目组主要看个人的兴趣。大部分社团都是需要缴纳团费的。工作室项目组都是小范围招新，需要你本身具备一定的学习能力。\\
	
	\subsection{大学可以做什么}
	很多同学刚进入学校并不知道这个组织是干什么的，只是觉得要做点什么，学长学姐又各种劝说加入相关组织的好处就加入了。
	又或者看着介绍很感兴趣，进入后发现跟介绍相比相差甚远，自己就是一个打杂的，想退出，又已经有了沉没成本。\\
	
	我的建议是先想想自己未来要做什么，再多向学长学姐咨询，多尝试才知道哪个是适合自己的。主动会为你赢得很多机会!\\
	
	如果你在大学期间不想做被水课耽误的咸鱼，这里有几条道路：
	
	\begin{itemize}
		\item 比赛。大部分比赛的知识认真学习三个月相关知识完全能获奖。尽量不要参与校级比赛（除非需要校赛优秀作品推荐），关注\href{https://tieba.baidu.com/p/9062077312?pn=1}{教育部认可的84项全国大学生学科竞赛名单}，来源学校官网。每年在国家级比赛获奖的同学都有很多。
		\item 赚钱。校内食堂有兼职岗位及勤工俭学的岗位。校内也有其他的赚钱渠道，一般是代理居多。（提示：慎防上当受骗。）
		\item 学习。继续学业，关注你的学科考研要靠什么，每年考研上岸的985 211的同学不在少数。想考公的学申论和行测，每年也有不少考公上岸的同学。
		\item 实习。关注岗位需求学习工作能力，提前进入公司实习。
	\end{itemize}
	
	
%	\begin{table}[!h]
%		\centering
	%	\caption{Sample table.} % 表格注释
	%	\begin{tabular}{cccc}
	%		\toprule
	%		Value 1 & Value 2 & Value 3 & Value 4\\
	%		\midrule
	%		odd     & odd   & odd & 1.00 \\
	%		even    & even  & even& 1.00 \\
	%		odd     & odd   & odd & 1.00 \\
	%		even    & even  & even& 1.00 \\
	%		\bottomrule
	%	\end{tabular}
%	\end{table}
		\newpage



\section{新生群常见问题解答} 
\textbf{Q}：学校男女比例是多少？

A：7：3

\textbf{Q}：学校有游泳馆吗？

A：没有

\textbf{Q}：有奖学金吗，奖学金类别、金额...？

A：普通本科生能获奖学金300-10000不等。贫困生还可申请助学金，2200-4400。

\textbf{Q}：学校地点在哪？

A：河北省张家口市桥东区南站街道朝阳西大街13号。

\textbf{Q}：学生就业率怎么样？

A：根据河北建筑工程学院2023年就业质量报告，毕业生的毕业去向落实率为 85.84$\%$

（包括：就业（签就业协议形式就业、签劳动合同形式就业、科研助
理、应征义务兵、其他录用形式就业、国家基层项目、地方基层项目、自由职业、自主创
业）、升学（国内升学、出国出境））。

\textbf{Q}：四级什么时候考，是强制的吗？

A：大一的同学下学期才可以开始报名，不是强制的，未通过也能顺利毕业。

\textbf{Q}：有早自习、晚自习吗？

A：没有早自习，部分院系有晚自习。

\textbf{Q}：能转专业吗？

A：根据学校《河北建筑工程学院优秀本(专)科生转专业实施细则(修订)》(校教[2017] 25号)规定，按转专业程序，可申请转入入学时对应的文理科专业，好像达到大一成绩在专业前百分之十就可以，具体见规定 或咨询学长学姐。

\textbf{Q}：书本订收费吗？

A：书本费自愿，跟学校走是200左右，可以自己网上订或者购买二手。



%\textbf{Q}：

%A：




		\newpage
\section{附录} 
	\subsection{各专业的对应院系}
		\begin{figure}[!h]
		\centering
		\includegraphics[width=0.25\textwidth]{mingdan.jpg}
		\caption{各专业的对应院系}
	\end{figure}
			\newpage
	
	\subsection{校园地图}
	\begin{figure}[!h]
	\centering
	\includegraphics[width=0.8\textwidth]{map.jpg}
	\caption{校园地图1}
\end{figure}

	\begin{figure}[!h]
	\centering
	\includegraphics[width=0.7\textwidth]{2024072555.jpg}
	\caption{校园地图2}
\end{figure}
			\newpage

	\subsection{物品携带}
	\frameboxbegin{报道携带} 
	 录取通知书、身份证、准考证、团组织关系证明、一寸两寸证件照（及电子），家庭情况调查表、贫困证明（录取通知书里面有，申请助学金用得到），生源地贷款回执单（针对贷款萌新后续贷款也需要）。
	\frameboxend
	\frameboxbegin{日常用品} 
	牙膏、牙刷、漱口杯，毛巾，洗衣液，挂钩，盆。充电器，充电宝，插座。床上用品。垃圾桶，垃圾袋。床边挂篮。床上桌。笔。雨伞。水杯。纸巾、湿巾。
	\frameboxend
	\frameboxbegin{关于校园卡} 
	校园卡办理免费。宿舍没wifi，本身的套餐流量不够用的话推荐办理校园卡，宿舍想用wifi需要充值网条也需要校园卡。提前问好套餐，在学长学姐那里办都有赠品。
	\frameboxend
\subsection{参与本文档建设}

本文档为学生自发创建，目的是帮助新一届同学更好的适应校园生活，如有侵权请联系文末删除。

本文档\href{https://github.com/ditian-tai/HBACE-Guide-for-Freshmen}{开源，地址https://github.com/ditian-tai/HBACE-Guide-for-Freshmen}，\textbf{转载请注明出处}。\\

现招募河北建筑工程学院新生指南文档建设者，你的帮助会让学弟学妹更快地融入校园生活，目前需要：
\begin{itemize}
	\item 文档美化，需要了解\LaTeX 的排版，使得文档排版整齐美观。
	\item 内容撰写，需要河北建筑工程学院大二及以上学生，有志于帮助新生，撰写内容需要保证真实，不得掺杂个人评价及情感。
	\item 图片提供，需要河北建筑工程学院大二及以上学生，有志于帮助新生，图片需要真实，美观，作为校园相关配图。
\end{itemize}


如有对文档内容的校正及建议也欢迎提出，“赠人玫瑰，手留余香”，欢迎更多的同学加入，你的名字及联系方式将出现在在本指南贡献者名单。联系方式：\href{mailto:cyselz@foxmail.com}{cyselz@foxmail.com}。
			\newpage
\section{本指南的主要贡献} 
\subsubsection*{关于大学的任何疑问均可进行咨询}
\subsubsection*{caytez，qq：1393883102（添加请备注姓名及专业）}
\subsubsection*{安哲，微信：ndgsyx862}

\end{document}
